\chapter{绪论}
\label{chap:intro}

\section{研究背景}

大语言模型驱动的智能体已从单轮问答走向\textbf{多步工具调用与环境交互}：在代码沙箱、终端、文件系统等真实环境中，智能体通过一连串"思考—调用工具—观察结果"的循环完成任务。要让智能体在这类交互式任务上不断变强，\textbf{强化学习}（RL）已成为主流手段——用交互轨迹上的奖励信号直接优化策略。

然而，把强化学习用于交互式智能体面临一个现实瓶颈：\textbf{它需要海量的交互轨迹与可信的奖励信号}。二者都不便宜。轨迹需要智能体在真实环境中反复试错地产生；奖励则需要判断"这条轨迹到底有没有把任务做好"。传统做法要么依赖\textbf{人工标注}奖励——昂贵、缓慢，无法支撑持续、大规模的策略更新；要么依赖\textbf{规则/自述}奖励——而一旦奖励只看智能体的自我陈述（"我已完成任务"），模型会迅速学会\textbf{谎报完成}这一典型的\textbf{奖励作弊}（reward hacking）行为：轨迹上写满"已成功创建文件"，而沙箱里根本没有那个文件。

因此，问题的核心变成：\textbf{能否让智能体在无人工标注的条件下，自主地产生经验、自主地获得可信奖励，从而自我进化？} 这正是本文所称\textbf{自进化强化学习系统}的目标——把"采样—奖励—更新"这一强化学习闭环做成一套无需人工介入、可持续运转的系统。

\section{研究问题}

本文聚焦：\textbf{如何构建一套面向交互式智能体的自进化强化学习系统，使其在无人工标注下自主产生多轮交互经验、并获得抗奖励作弊的可信奖励，进而以标准强化学习稳定地更新策略？}

这一总问题可拆为三个子问题：（1）如何组织采样，让智能体在真实环境中自主产生\textbf{有信息量的多轮交互经验}；（2）如何构造奖励，使其\textbf{锚定环境真实状态}而非智能体自述，从机制上抑制奖励作弊；（3）如何将上述采样与奖励机制整合为一套可稳定运转的强化学习系统。

\section{主要贡献}

\begin{enumerate}
  \item \textbf{多智能体采样架构}（第\ref{chap:method}章）：执行者在代码沙箱中以 ReAct 方式多步调用工具；观察者以沙箱执行前后的\textbf{状态差分}为客观依据生成状态报告；提问者扮演具人设的模拟用户驱动多轮追问。三者协同把"一次采样"扩展为"一段自主多轮交互经验"。
  \item \textbf{差分驱动的可信奖励}（第\ref{chap:method}章）：奖励由外部冻结大模型裁判依据观察者提供的\textbf{环境真实状态差分}打分，四维乘性聚合、安全二元门控。以真实状态而非自述为锚点，从机制上抑制奖励作弊。
  \item \textbf{系统实现与端到端验证}（第\ref{chap:system}、\ref{chap:exp}章）：将上述采样与奖励机制实现为一套完整的智能体强化学习系统，推理与训练异步分离，并在 Qwen3.5-9B 上端到端验证整套自进化闭环能够稳定运转。
\end{enumerate}

\section{论文组织}

第\ref{chap:related}章综述相关工作；第\ref{chap:problem}章形式化问题；第\ref{chap:method}章详述方法；第\ref{chap:system}章介绍系统实现；第\ref{chap:exp}章给出实验设计与验证；第\ref{chap:conclusion}章总结展望。
